\section{Conclusion}
\label{sec:conclusion}


We have introduced \ourmethod{}, a unified feedforward model for 4D reconstruction from a pair of unposed images. By predicting a holistic dynamic 3D Gaussian Splatting representation, \ourmethod{} effectively leverages self-supervision to overcome challenges like data scarcity and occlusions. \ourmethod{} significantly outperforms prior work in estimating 3D geometry and motion, while uniquely enabling high-fidelity 4D spatio-temporal interpolation. Our work demonstrates the power of unified, explicit representations for dynamic scene understanding, opening promising avenues for extending this framework to long-sequence videos and more complex motion models.


\paragraph{Future work.}
We identify several directions for future explorations. A naive extension to long sequences would be memory-intensive due to the linear growth of Gaussian Splats; future work could explore a compact scene representation~\citep{An:2025:C3G,Xu:2025:RSL}. 
Also, our model's assumption of linear motion and constant brightness is best suited for short time intervals. This could be addressed by incorporating learnable, non-linear motion models and time-varying Gaussian attributes to handle more complex dynamics and photometric changes.
Moreover, rendered images, motion, and geometry at novel view and time can be included as additional supervision objectives to enhance spatio-temporal consistency, when their annotations (\eg from multi-view video datasets) are available. 

\subsubsection*{Acknowledgments}
We sincerely thank Richard Szeliski, Junyi Zhang, Michael Rubinstein, Richard Tucker, and Linyi Jin for their helpful discussions and technical assistance throughout the project. We also appreciate the anonymous reviewers for their constructive feedback during the review process.